\begin{center}
  \large\textbf{ABSTRAK}
\end{center}

\addcontentsline{toc}{chapter}{ABSTRAK}

\vspace{2ex}

\begingroup
% Menghilangkan padding
\setlength{\tabcolsep}{0pt}

\noindent
\begin{tabularx}{\textwidth}{l >{\centering}m{2em} X}
  Nama Mahasiswa    & : & \name{}         \\

  Judul Tugas Akhir & : & \tatitle{}      \\

  Pembimbing        & : & 1. \advisor{}   \\
                    &   & 2. \coadvisor{} \\
\end{tabularx}
\endgroup

Indonesia mengalami frekuensi gempa bumi tinggi yang menyebabkan perubahan
struktural signifikan pada lingkungan \emph{indoor} bertingkat. Robot otonom
dapat berperan menavigasi area \emph{indoor} pasca gempa yang tidak aman
bagi manusia. Untuk melakukan misi tersebut, robot memerlukan kemampuan
navigasi yang mampu bekerja pada lingkungan pasca gempa yang bertingkat.
Tugas akhir ini membahas pengembangan metode navigasi menggunakan \emph{layered
costmap} supaya robot dapat melintasi lingkungan tersebut. Skema ini
merepresentasikan lingkungan bertingkat sebagai kumpulan peta \emph{region}
berbentuk \emph{occupancy grid} yang digabung menjadi satu \emph{map}. Dalam
\emph{costmap} berlapis, tiap lapisan pada \emph{costmap} menyimpan informasi
berbeda seperti rintangan statis, rintangan dinamis, lapisan inflasi, dan
lapisan transisi antar \emph{map}, yang memungkinkan navigasi antar
\emph{region}. Dalam arsitektur navigasi yang dikembangkan, setiap lantai pada
lingkungan \emph{indoor} direpresentasikan oleh satu \emph{region} individual.
Pendekatan ini menyederhanakan informasi lingkungan 3D menjadi representasi 2D
yang dapat dikelola, sehingga memungkinkan penggunaan sensor \emph{LiDAR} 2D
untuk menentukan lingkungan robot. Algoritma navigasi kemudian mengintegrasikan
\emph{map} 2D tiap \emph{region} ini menjadi satu kesatuan dengan menggunakan
\emph{costmap} transisi, sehingga memungkinkan robot untuk memahami dan
menavigasi lingkungan multi-lantai. Sistem ini memungkinkan robot untuk
menavigasi dari satu \emph{region} ke \emph{region} lain dengan (1) estimasi
pose menggunakan \emph{iterative closest point} (ICP) dan \emph{region
switcher}, (2) \emph{path planning} yang membuat rute antar \emph{region}
menggunakan \emph{path router}, dan (3) \emph{path tracking} yang menggunakan
\emph{dynamic window approach} (DWA) \emph{planner} untuk menghasilkan
\emph{trajectory} yang diikuti robot.
Dari hasil pengujian pada simulator MuJoCo, robot berhasil melakukan navigasi
dari \emph{region} A sampai C dengan tingkat keberhasilan 80\% (8 dari 10
percobaan), waktu tempuh rata-rata 199.35 detik, dan waktu terbaik 141.01
detik.

% % Ubah kata-kata berikut dengan kata kunci dari tugas akhir
Kata Kunci:  \emph{Navigasi, Costmap, Robot, Antar lantai}
